\section*{Abstract}

London Fire Brigade publishes detailed incident and mobilisation records, but the public data are difficult to interpret as evidence for response-performance questions: borough demand, incident type, response-time denominators, appliance-level mobilisation, and spatial uncertainty all interact. This project asks how an uncertainty-aware visual analytics dashboard can support structured inspection of London Fire Brigade incident demand and first-pump response performance using public open data.

The project builds a reproducible Python data pipeline over the 2024--2026 LFB Incident Records, LFB Mobilisation Records, London borough boundaries, and a public fire-station address/postcode table. The processed artefacts feed a Flask API and a D3.js linked-view dashboard with borough map, monthly trend, response-time distribution, and ranking-table panels. The dashboard is deliberately framed as an exploratory evidence surface rather than an operational dispatch, routing, forecasting, optimisation, or station-relocation system. Data-quality constraints are carried through the pipeline and report: only 37.9\% of incident rows have precise coordinates, mobilisation-derived claims are filtered to matched canonical incidents, and the station layer distinguishes incident-derived assignment footprints from postcode-geocoded public station-address proximity.

The dashboard exposes borough-level response-performance disparities and the limits attached to them. Among 278,468 incidents with a recorded first-pump time, the author's share above 360 seconds is 32.4\% (95\% Wilson CI 32.3--32.6\%); this is not LFB's official per-incident failure rate. The highest borough values are Hillingdon (52.2\%), Bromley (47.6\%), and Richmond upon Thames (45.2\%). The ranking remains stable when false alarms are excluded. The dashboard shows recorded denominators, confidence intervals, small-sample warnings, and the partial final month. Coordinate missingness is kept with point-level diagnostics rather than used to qualify borough aggregates. Mobilisation and historical station-address evidence provide separate context. The evaluation supports a narrow claim about linked comparison; several original tasks exceeded the implemented interface, and raw per-session notes are not included in the submitted workspace.

\clearpage
